\chapter{Ownership and Borrowing}

The way a programming language manages memory fundamentally defines its character, performance, and safety guarantees. In languages like C and C++, memory management is manual: the programmer must explicitly request and release memory. While this offers maximum control, it inevitably leads to complex bugs such as \emph{dangling pointers} (using memory after it has been freed), \emph{double frees} (freeing the same memory twice), and memory leaks. On the other end of the spectrum, languages like Java or Python use a Garbage Collector (GC), which automatically tracks and frees unused memory at runtime. This removes the burden from the programmer but introduces performance overhead and unpredictable pauses.

Rust introduces a paradigm-shifting alternative: \textbf{Ownership}. By shifting memory management checks to compile-time via a component known as the \emph{Borrow Checker}, Rust achieves the performance of manual memory management without its inherent dangers. The programmer is forced to explicitly consider the lifespan and ownership of every value during the writing phase. While this may initially feel like a ``difficulty gradient,'' it acts as a massive safety net. As the professor emphasized: ownership, borrowing, and lifetimes are concepts that intrinsically belong to programming---they exist in every language---but Rust is the first to make them \emph{explicit}. Other languages allow the programmer to describe things that break at runtime because they never force this awareness. Rust's compiler \emph{rejects} programs that do not satisfy these rules.

\begin{insightbox}
\textbf{The Red Light Analogy} \\
The professor compared ignoring ownership rules to running a red light: \emph{``It doesn't automatically mean getting hit by a truck. Sometimes you get lucky. And so you develop the false awareness that your program works---but in reality, it is intrinsically broken. It will break at the worst possible moment.''}  In Rust, the code simply will not compile if memory safety rules are violated, preventing you from ever reaching that dangerous intersection.
\end{insightbox}

\begin{insightbox}
\textbf{The Hypercorrectivity Analogy} \\
The professor explained why bugs escape programmers' attention with a linguistic analogy: \emph{``We are naturally hypercorrective. When we read a text, we give it a quick look, intuit what should be written, and if it mentally `wires' correctly, we are happy. People don't notice missing accents or misplaced apostrophes---their brains auto-correct. The same happens with code: we glance at it, we intuit that it should work, and we move on. We don't have the mental strength to pay attention to every detail.''} The Borrow Checker, unlike the human mind, does not suffer from hypercorrectivity---it meticulously checks every single path.
\end{insightbox}

% ================================================================
\section{The Core Concept of Ownership}
% ================================================================

The foundation of Rust's memory model rests on three deceptively simple rules:
\begin{enumerate}
    \item Every value in Rust has a variable that is called its \emph{owner}.
    \item There can only be one owner at a time.
    \item When the owner goes out of scope, the value will be dropped (its memory and resources are freed).
\end{enumerate}

Note the choice of subject and object: \emph{the value is possessed by a variable}. If a value exists, it must reside somewhere in memory---a value without representation ontologically cannot exist. The memory in which that value resides \emph{is} the variable. A particular component of the compiler, the \textbf{Borrow Checker}, verifies this at every point of your program during compilation. It reads your code with one specific objective: is there any point where the programmer behaves as if a value is not possessed by exactly one variable? If it finds such a point, it blocks compilation and explains exactly why.

\subsection{What Does ``Possessing a Value'' Mean?}

Possessing a value means being \textbf{responsible for its release}. The professor made this crucial distinction: the problem is not what you do with a value while you are using it---you do whatever you want. The problem is: \emph{when you no longer need it, you must ensure that the memory (and any associated resources) returns to its legitimate owner}.

Sometimes the cleanup goes beyond simply releasing memory. Consider opening a file: the operating system builds internal data structures---tables that map the file handler to disk locations, buffers for I/O operations, and so on. The value you hold is just a small number (a file handler). Beyond releasing the memory where that small number lives, you must also ensure the operating system cleans up all the information it is keeping associated with that handler. The file must be \emph{closed}. The same principle applies to network sockets, database connections, and any other system resource.

\begin{insightbox}
\textbf{The House Ownership Analogy} \\
During a Q\&A session, the professor illustrated ownership with a vivid analogy: \emph{``The house in which you live---whose owner is it? The landlord's. But you can live in it. Why? Because you have a contract. You have permission to be there. You don't own the house---you are borrowing it.''} This maps directly to Rust's reference system: variables can grant temporary access (references) to their values without transferring ownership. But the owner remains the one who is responsible for the ``house'' when the ``tenant'' leaves.
\end{insightbox}

\begin{insightbox}
\textbf{The Differentiated Waste Collection Analogy} \\
Possessing a value in Rust means being responsible for its disposal. Think of it like differentiated waste collection: \emph{``What do I do with paper? I put it in the yellow bin. What do I do with glass? In the blue bin. What about electronics, or hazardous materials like paint residues? Those require a trip to the eco-center.''} Based on the typology of what must be disposed of, different operations are required. Similarly, some data types have trivially simple disposal (just pop the stack), while others require complex cleanup procedures. When we introduce a type in Rust, we can declare specialized disposal instructions by implementing the \texttt{Drop} trait.
\end{insightbox}

\subsection{When Does Release Happen?}

The release of a value occurs in three situations:

\begin{enumerate}
    \item \textbf{End of scope:} When the variable goes out of scope---that is, when program execution reaches the closing brace of the block in which the variable was declared---the value is released.
    \item \textbf{Reassignment:} When a variable is assigned a new value while still in scope, the \emph{previous} value must be cleaned up first. Only then is the new value assigned. Failing to clean up the previous value would result in a resource leak.
    \item \textbf{Explicit drop:} You can accelerate the release by calling \texttt{drop(variable)}. This tells Rust: ``from a syntactic point of view, this variable would still be usable for the next 25 lines, but I no longer need it---drop it now.'' After the call, Rust marks the variable as unusable for the remainder of its scope.
\end{enumerate}

\begin{warningbox}
\textbf{Why release matters} \\
Resources occupied by a variable must become available to others. If a variable occupies memory, that memory could be reused by others doing something entirely different. If a file handle is not released, the operating system accumulates garbage in its internal tables. Release is not optional---it is \emph{mandatory} for correctness.
\end{warningbox}

% ================================================================
\section{Ownership in Action: The Vec Example}
% ================================================================

Let us trace a complete ownership lifecycle through the professor's detailed example of a vector.

\subsection{Vec::new() vs.\ Vec::with\_capacity()}

Before diving into the example, the professor distinguished two ways to construct a vector:

\begin{itemize}
    \item \textbf{\texttt{Vec::new()}} --- creates an empty vector with \emph{no} initial heap allocation. It starts with capacity 0 and length 0. No memory is requested from the operating system until the first element is pushed.
    \item \textbf{\texttt{Vec::with\_capacity(n)}} --- creates an empty vector but \emph{pre-allocates} space for \texttt{n} elements on the heap. It starts with capacity \texttt{n} and length 0. The heap block is ready, even though no elements have been inserted yet.
\end{itemize}

The professor chose \texttt{with\_capacity(4)} for the example to illustrate the reallocation mechanism clearly.

\begin{lstlisting}[language=Rust]
{
    let mut v = Vec::with_capacity(4);
    for i in 1..=5 {
        v.push(i);
    }
    println!("{:?}", v);
} // 'v' goes out of scope here
\end{lstlisting}

\textbf{Step-by-step memory walkthrough:}

\begin{memoryblock}
\textbf{Step 1 --- Declaration: \texttt{let mut v = Vec::with\_capacity(4);}} \\
The stack grows by 24 bytes to accommodate the variable \texttt{v}. The \texttt{Vec} structure contains three 64-bit fields on the stack:
\begin{itemize}
    \item An 8-byte \textbf{pointer} to a heap-allocated buffer
    \item An 8-byte \textbf{capacity} field (currently = 4)
    \item An 8-byte \textbf{length} field (currently = 0, the vector is empty)
\end{itemize}
Because we used \texttt{with\_capacity(4)}, the operating system has already been asked to provide a block of 4 slots on the heap. At this moment, \texttt{v} is the \emph{possessor} of this composite structure: a bit of stack and a bit of heap.

\textbf{Step 2 --- Pushing elements 1 through 4:} \\
Each \texttt{v.push(i)} places the value \texttt{i} into the next available heap slot and increments the length counter. After four pushes: capacity = 4, length = 4. All slots are occupied.

\textbf{Step 3 --- Pushing element 5 (reallocation):} \\
When we try \texttt{v.push(5)}, the vector discovers it has no more room. It \emph{doubles} its capacity: it asks the operating system for a new block of 8 slots, copies the existing 4 elements into the new block, returns the old block of 4 to the OS, then inserts the 5th element. The stack pointer is updated to point to the new block. Capacity becomes 8, length becomes 5. If further pushes filled slots 6, 7, 8, the next push would trigger another doubling to 16 slots, and so on.

\textbf{Step 4 --- End of scope (the closing brace):} \\
The stack needs to contract by 24 bytes. But \texttt{v} is not a simple variable---it has data on the heap. The \texttt{Vec} type is labeled with the \texttt{Drop} trait: there are ``dismantling instructions on its box.'' Rust reads the type, sees the drop label, and \emph{before} contracting the stack, executes the specialized drop procedure. The drop logic checks: if the capacity is non-zero, then the internal pointer is valid and the heap block must be freed. In this case, capacity is 8, so 8 slots worth of heap memory are returned to the OS. Only then does the stack contract.
\end{memoryblock}

% AI-IMAGE-REQUEST: Detailed diagram showing a Vec's memory layout across stack and heap. The stack portion shows three 8-byte fields (pointer, capacity=8, length=5). An arrow from the pointer field leads to a heap buffer with 8 slots, 5 of which contain values 1-5 and 3 are empty. Show the reallocation process: old block of 4 being copied to new block of 8, with the old block crossed out and returned to OS.
\begin{figure}[!ht]
    \centering
    \includegraphics[width=0.85\textwidth]{images/ch04_vec_memory_layout_and_reallocation.png}
    \caption{Memory layout of a \texttt{Vec}: stack metadata (pointer, capacity, length) and heap buffer, including the reallocation process when capacity is exceeded.}
\end{figure}

\begin{exambox}
\textbf{Understanding the Drop Trait for Vec} \\
The compiler knows what to do because the \texttt{Vec} type is \emph{decorated} with drop instructions. The compiler knows the type of every variable and, consequently, when it arrives at the point where a \texttt{Vec} must be released, it applies the appropriate rules. The drop trait for \texttt{Vec} checks: ``if my capacity is different from zero, the pointer in my belly is valid and must be freed.'' This is the essence of \textbf{RAII} (Resource Acquisition Is Initialization) in Rust.
\end{exambox}

% ================================================================
\section{Transferring Ownership: Moves}
% ================================================================

Because a value can have only one owner, assigning an existing variable to a new one fundamentally changes the state of the program. This is the concept of \textbf{move}.

\subsection{The Three Forms of Assignment}

The professor identified three syntactic forms of assignment, all operationally identical:

\begin{enumerate}
    \item \textbf{Explicit assignment:} \texttt{let b = a;} --- the value moves from \texttt{a} to \texttt{b}.
    \item \textbf{Passing a parameter:} \texttt{f(a)} --- the value is copied from \texttt{a}'s position to the function's parameter slot on the stack.
    \item \textbf{Returning a result:} \texttt{return a;} --- the value is passed from the function's local scope to the caller.
\end{enumerate}

In all three cases, when the type is not \texttt{Copy}, the \emph{possession transfers} from the source to the destination. The source becomes inaccessible. The compiler tracks this: it will do nothing special when the source reaches its end of life (no drop instructions), and will instead execute the specialized drop on the \emph{destination} when the destination's life ends.

The professor observed an asymmetry in how surprising these moves feel: \emph{``The move on return is not surprising---we don't even get surprised that the function disappears. We get more surprised on the parameter. If I call \texttt{f(a)}, on the next line I can't read \texttt{a} anymore.''} This surprise fades with experience, but it is the most common stumbling block for new Rust programmers.

\subsection{Detailed String Move Example}

\begin{lstlisting}[language=Rust]
let mut s1 = "Hello".to_string(); // s1 owns the string
println!("{}", s1);               // Works fine
let s2 = s1;                      // Ownership MOVES to s2
// println!("{}", s1);            // COMPILE ERROR: value used after move
println!("{}", s2);               // Works fine: s2 is the new owner
\end{lstlisting}

\begin{memoryblock}
\textbf{Step-by-step memory trace:}

\textbf{1.} \texttt{let mut s1 = "Hello".to\_string();} --- The stack grows by 24 bytes for \texttt{s1}. A \texttt{String}'s internal structure is very similar to a \texttt{Vec}: a pointer to the heap, a capacity, and a length. The \texttt{to\_string()} method allocates exactly 5 bytes on the heap (``Hello'' has 5 characters), initializes them, and sets capacity = 5, length = 5.

\textbf{2.} \texttt{let s2 = s1;} --- The stack grows by another 24 bytes for \texttt{s2}. What goes into these 24 bytes? \emph{A bitwise copy} of everything that was in \texttt{s1}---the pointer, the capacity, the length. This is a ``dumb copy'' (\emph{copia becera}), and it is extremely cheap: just a \texttt{memcpy} of 24 bytes.

\textbf{3.} Without the move rule, we would be in an extremely uncomfortable situation: two variables pointing to the same heap block. Who releases it? If both try, it is a double free. If neither does, it is a leak. The solution: the compiler marks \texttt{s1} as \textbf{inaccessible} from this line forward. Any attempt to read \texttt{s1} results in a compile error.

\textbf{4.} When \texttt{s2} reaches the end of its scope, the compiler knows it is the owner: it executes the \texttt{Drop} procedure (freeing the heap block), then contracts the stack. When \texttt{s1} reaches the end of scope, the compiler \emph{omits} the specialized drop---\texttt{s1} still physically exists on the stack (the stack must still contract by 24 bytes), but it no longer ``has a value'' in the ownership sense. It is inaccessible.
\end{memoryblock}

% AI-IMAGE-REQUEST: Three-panel diagram showing: (1) s1 on stack with arrow to "Hello" on heap, (2) after move: s1 grayed out/crossed, s2 on stack with arrow to same "Hello" on heap, (3) at end of scope: s2 triggers drop, heap freed, stack contracts. s1's drop is skipped.
\begin{figure}[!ht]
    \centering
    \includegraphics[width=\textwidth]{images/ch04_string_move_ownership_transfer.png}
    \caption{Ownership transfer of a \texttt{String}: the bitwise copy of stack data, invalidation of the source, and how drop is executed only by the new owner.}
\end{figure}

\begin{insightbox}
\textbf{How this resolves C's pointer ambiguity} \\
The professor emphasized that this mechanism removes one of the essential ambiguities of C pointers: \emph{``We said pointers are ambiguous---we don't know if they are valid, if we have to release them, or if someone else already released them. Here the compiler takes awareness and tells us: `Calm down, I'll track it.' One and only one pointer is the possessor of that block. The other one can't touch it anymore.''}
\end{insightbox}

\begin{insightbox}
\textbf{How this prevents double frees} \\
The professor explicitly connected single ownership to double-free prevention: \emph{``Since ownership belongs to one and only one variable, and only the owner releases, there is no possibility of a double release---which is another of the problems of C. Because if I have a pointer, a copy of that pointer, and someone has already freed it, and I don't know this, I could call \texttt{free} a second time. But if we call \texttt{free} twice, in the meantime that memory has been given to others---I'll break everything.''} In Rust, the compiler statically guarantees that exactly one entity performs the release, exactly once.
\end{insightbox}

\subsection{Reassignment After Move}

A moved variable is not ``dead forever.'' If you assign a \emph{new} value to it, it comes back to life:

\begin{lstlisting}[language=Rust]
let mut s1 = "Hello".to_string();
let s2 = s1;              // s1 is moved, now inaccessible

// Read access to s1 here would fail:
// println!("{}", s1);    // ERROR: value used after move

s1 = "World".to_string(); // s1 is REBORN with a new value
println!("{}", s1);        // Works: prints "World"
println!("{}", s2);        // Works: prints "Hello"
\end{lstlisting}

After the reassignment, \texttt{s1} has a completely new value---a fresh heap allocation containing ``World''. The old ``Hello'' heap block belongs entirely to \texttt{s2}. When both variables reach the end of scope, \texttt{s2} will free the ``Hello'' block and \texttt{s1} will free the ``World'' block. Each pointer releases its own block, exactly once.

The professor noted: \emph{``After the move, read access to the source will cause a compilation error, but write access---assigning a new value---will succeed and result in the variable acquiring ownership of the new thing.''}

The professor also clarified an important detail about the heap: \emph{``When you reassign \texttt{s1} with a new value, the old heap block is not reused. A completely new arrow is born, pointing to a different location on the heap. The only way to modify the variable is to use one of the creation functions of the type, which paints all the data structure---the three fields of the String---with something valid. The old block was separated, so it won't bother us.''}

\subsection{Shadowing vs.\ Reassignment}

The professor carefully distinguished between \texttt{let s1 = s2;} (shadowing with \texttt{let}) and \texttt{s1 = s2;} (reassignment without \texttt{let}):

\begin{itemize}
    \item \textbf{\texttt{let s1 = s2;}} introduces a \emph{new} variable with the same name. The old \texttt{s1} still physically exists on the stack but becomes inaccessible---not because it lost its value via a move, but because it is \emph{hidden} (\emph{shadowed}) by the new declaration. The new \texttt{s1} occupies additional space further on the stack.
    \item \textbf{\texttt{s1 = s2;}} (without \texttt{let}) performs the \emph{reverse movement}: the value moves back from \texttt{s2} into the existing \texttt{s1} variable. Now \texttt{s2} loses possession and becomes inaccessible, while \texttt{s1} regains ownership.
\end{itemize}

\begin{warningbox}
\textbf{Shadowing creates a new stack slot} \\
When you write \texttt{let s1 = s2;} inside the same scope, you are not reusing \texttt{s1}'s stack slot. The new \texttt{s1} occupies \emph{additional} space on the stack. The old \texttt{s1} is still there, hidden. When the block ends, the stack contracts and the old one becomes syntactically visible again (if the block was an inner block). This is a subtle difference from simple reassignment.
\end{warningbox}

\subsection{Implicit Moves: The Operator Trap}

Moves are not limited to direct assignments. Any operation that syntactically looks like an operator may actually be a function call that takes ownership:

\begin{lstlisting}[language=Rust]
let s = String::from("Hello");
let b = s + " world"; // IMPLICIT call to concat function!
// s is NO LONGER accessible here - it was moved into concat
println!("{}", b);     // Works: prints "Hello world"
// println!("{}", s);  // COMPILE ERROR: s was moved
\end{lstlisting}

The \texttt{+} operator on strings is actually an implicit call to a concatenation function (internally \texttt{add}/\texttt{concat}). Because you are passing \texttt{s} as a parameter, ownership transfers to the function. From that point on, \texttt{s} is inaccessible---the heap block it was using may have been reused by the result \texttt{b}.

The professor traced the chain of possession explicitly: \emph{``The possession of that block has potentially passed into the hands of \texttt{b}---that is, passed into the hands of \texttt{concat}, which could have been transferred to \texttt{b} in the last analysis. So \texttt{s} has no right to look inside anymore.''} Inside \texttt{s} the pointer still physically exists on the stack, but it is \emph{inaccessible} because the heap memory it was pointing to may have been reused by the result.

Contrast this with numeric types:

\begin{lstlisting}[language=Rust]
let i = 42;
let j = i + 1; // Also an implicit function call (add)
println!("{}", i); // Works! i is still 42
println!("{}", j); // Works! j is 43
\end{lstlisting}

The difference? Integers implement the \texttt{Copy} trait. They have no special disposal instructions, no heap allocations, no ambiguity. Both \texttt{i} and \texttt{j} remain independently usable.

\begin{examprioritybox}
\textbf{Implicit Moves and Reassignment} \\
Moves occur in three ways: direct assignment (\texttt{let s2 = s1;}), passing arguments to a function (\texttt{f(s1)}), and returning a value from a function (\texttt{return s1;}). Furthermore, implicit function calls through operators (e.g., \texttt{let b = s + " world";}) also trigger moves. Once you pass a non-copyable variable to a function, you yield ownership to that function's parameter, rendering the original variable inaccessible in the calling scope.
\begin{itemize}
    \item \textbf{Common Pitfall:} Students often assume variables remain usable after passing them to a function or using them in an operator expression. In Rust, this causes a ``value used after move'' compile error.
    \item \textbf{Expected Mastery Level:} Very High. You must statically trace the ownership of every value. Recognize that \texttt{Copy} and \texttt{Drop} traits are mutually exclusive. Copy types are the \emph{exception}, not the rule.
\end{itemize}
\end{examprioritybox}

% ================================================================
\section{The Trait System: Drop, Copy, and Clone}
% ================================================================

Types in Rust carry ``labels'' that describe how they should be treated. These labels are called \textbf{traits}. Three traits are central to the ownership system:

\subsection{The Drop Trait}

The \texttt{Drop} trait indicates that a type requires \textbf{specialized cleanup} when it is released. For a \texttt{Vec}, this means freeing its heap buffer. For a \texttt{File}, this means closing the file handle and releasing OS resources. When a variable of a \texttt{Drop} type goes out of scope, Rust automatically calls its \texttt{drop} method before contracting the stack.

The professor noted a naming subtlety: the \texttt{Drop} \emph{trait} (capital D) is the label/interface; the \texttt{drop()} \emph{method} (lowercase d) is the actual destruction operation that gets called. They are related but distinct---the trait declares the capability, the method implements it.

\subsection{The Copy Trait}

The \texttt{Copy} trait indicates that a type can be \textbf{freely duplicated}---when an assignment is made, the original remains accessible. Why? Because the value has no special meaning beyond its raw bit representation. There are no heap pointers to manage, no OS resources to clean up. The copy is harmless.

The professor explained: \emph{``The numeric types---\texttt{i32}, \texttt{f64}, understood as integers and floats on which I do multiplications and divisions---can be easily copied. There are no particular dismantling instructions. They are among those things I can toss directly into the regular trash bin. No differentiated collection needed.''}

The professor further clarified why strings cannot be \texttt{Copy}: \emph{``The copy of the pointer bothers me because that pointer, even though it is a bit sequence like all other values, has a particular meaning---it tells me something about the heap. It possesses the block to which it points. When should I destroy the heap block? When \texttt{s1} ends its life, or when \texttt{s2} ends its life? Both seem to have the right to that thing, and it's not okay. So the String is not copyable---it's only movable.''}

\subsection{The Mutual Exclusivity of Drop and Copy}

\begin{warningbox}
\textbf{Drop and Copy are mutually exclusive} \\
If a type implements \texttt{Drop}, it \textbf{cannot} implement \texttt{Copy}. If it implements \texttt{Copy}, it \textbf{cannot} implement \texttt{Drop}. The compiler enforces this constraint. The logic is straightforward: if a type needs specialized cleanup (Drop), then having two independent copies creates ambiguity---which copy is responsible for the cleanup? If a type is trivially copyable (Copy), it by definition requires no special cleanup.
\end{warningbox}

There is also a third category: types that are \textbf{neither Drop nor Copy}. These types have no specialized cleanup instructions, but they are not freely copyable either---for other semantic reasons. Assignment of these types still results in a move.

\begin{center}
\begin{tabular}{|l|c|c|c|}
\hline
\textbf{Type Category} & \textbf{Drop?} & \textbf{Copy?} & \textbf{On Assignment} \\
\hline
\texttt{i32}, \texttt{f64}, \texttt{bool}, \texttt{char} & No & Yes & Bitwise copy (original accessible) \\
\texttt{String}, \texttt{Vec<T>}, \texttt{File} & Yes & No & Move (original inaccessible) \\
Neither Drop nor Copy & No & No & Move (original inaccessible) \\
\hline
\end{tabular}
\end{center}

\subsection{Who Decides the Trait Labels?}

The professor stressed an important point: \emph{``A type is not Copy because the compiler decides---it is because the programmer labels it. The compiler merely verifies that the label is compatible.''} When using existing standard library types (\texttt{String}, \texttt{Vec}, \texttt{i32}, \texttt{File}, etc.), the labels are already set by the library authors. But when you introduce your own types---which happens starting from the next week of the course---it will be \emph{you} who decides ``what label to use,'' and you must do it knowing what you are putting inside the type. The compiler enforces consistency (e.g., you cannot label a type both \texttt{Drop} and \texttt{Copy}), but the initial decision is the programmer's.

\subsection{Copy Propagation in Composite Types}

The \texttt{Copy} property propagates through composite types:

\begin{itemize}
    \item \textbf{Tuples:} A tuple is \texttt{Copy} if \emph{all} its elements are \texttt{Copy}. For example, \texttt{(i32, f64)} is Copy because both components are Copy. But \texttt{(i32, String)} is \emph{not} Copy because \texttt{String} is not Copy.
    \item \textbf{Arrays:} A fixed-size array \texttt{[T; N]} is \texttt{Copy} if \texttt{T} is \texttt{Copy}. Arrays are homogeneous and fixed-size (known at compile time).
    \item \textbf{Tuples vs.\ Arrays:} Tuples can be heterogeneous (\texttt{(i32, f64, bool)}), arrays are homogeneous. Both propagate the Copy trait from their elements.
\end{itemize}

The professor gave the example of 3D coordinates: \emph{``If I create a triple of numbers to represent a point in space, the numbers are Copy, and consequently the triple representing the three coordinates x, y, z is automatically Copy.''}

The professor also introduced the general term: \emph{``A tuple is the generic name to indicate a pair, a triple, a quadruple, an n-tuple. A triple always has three things. Those three things can be potentially heterogeneous.''} This heterogeneity is what distinguishes tuples from arrays---arrays are always homogeneous.

\subsection{What ``Move'' and ``Copy'' Actually Do at the Machine Level}

At the hardware level, both a move and a copy are \emph{identical}: a \texttt{memcpy} of N bytes from source to destination. The only difference is a \textbf{compiler annotation}: when a type implements \texttt{Copy}, the compiler does \emph{not} add the ``data is now inaccessible'' mark to the source. That is the entire distinction. No different machine instructions, no runtime cost difference---just a compile-time bookkeeping difference.

\subsection{The Clone Trait: Explicit Deep Copy}

Sometimes you genuinely need an independent copy of a non-Copy type. The \texttt{Clone} trait provides this through an explicit method call:

\begin{lstlisting}[language=Rust]
let s1 = "Hi".to_string();
let s2 = s1.clone(); // Deep copy: new heap allocation
// Both s1 and s2 are independently usable
println!("s1 = {}, s2 = {}", s1, s2);
\end{lstlisting}

\begin{memoryblock}
\textbf{Clone vs.\ Move in memory:}

When we wrote \texttt{let s2 = s1;} (move), only the 24 stack bytes were copied. The heap block remained singular---and the source was invalidated.

When we write \texttt{let s2 = s1.clone();}, the clone operation:
\begin{enumerate}
    \item Allocates a \emph{new} block on the heap, as large as the original
    \item Copies all the characters from the original heap block to the new one
    \item Stores the pointer to the new block in \texttt{s2}, along with the appropriate capacity and length
\end{enumerate}
Now \texttt{s1} and \texttt{s2} each have their own independent heap block. Modifying one does not affect the other.
\end{memoryblock}

\begin{lstlisting}[language=Rust]
let mut s1 = "Hi".to_string();
let s2 = s1.clone();       // s2 gets independent copy
s1.push('!');               // Modifies only s1's heap data
println!("s1 = {}", s1);   // "Hi!"
println!("s2 = {}", s2);   // "Hi" (unchanged)
\end{lstlisting}

The professor traced the push operation in detail: when \texttt{s1.push('!')} is called, \texttt{s1} only has a 2-byte heap block (for `H' and `i'), fully occupied. Since \texttt{s1} is mutable, it requests a larger block from the OS (e.g., 4 bytes), copies the existing 2 characters, adds the exclamation mark in the third slot, and frees the old 2-byte block. Meanwhile, \texttt{s2}'s heap block remains untouched. The professor concluded: \emph{``So when I clone, I'm separating the objects. When I had the movement, instead, they were the same object with one of the paths disabled.''}

\begin{warningbox}
\textbf{Clone can be expensive!} \\
The professor used a vivid example: \emph{``While a string move always copies exactly 24 bytes (fixed cost), cloning a string depends on its length. If the string contained the entire Divine Comedy, the clone would require allocating a huge heap block and copying every single byte. The cost is proportional to the data size.''}  Always consider performance implications when using \texttt{clone()}.
\end{warningbox}

\subsection{The Relationship Between Clone and Copy}

There is a logical dependency:
\begin{itemize}
    \item Only \texttt{Clone}-able types can possibly be \texttt{Copy}. If a type cannot be cloned, it means its semantics fundamentally exclude replication---so there is no reason it could be copyable.
    \item When a type is both \texttt{Clone} and \texttt{Copy}, their implementations are generally identical (but this is not mandatory).
    \item \texttt{Copy} is used implicitly on assignment/passing. \texttt{Clone} requires an \emph{explicit} \texttt{.clone()} call.
    \item Non-cloneable types are relatively rare but exist for specific semantic reasons.
\end{itemize}

The professor added: \emph{``It is not obligatory to be cloneable, just as it is not obligatory to have particular recycling rules, and just as it is not obligatory to be copyable. A good portion of the types are cloneable, and that's fine. Some are not---for serious reasons which we will understand better later.''}

The professor stressed: \emph{``Copy types are the exception, not the rule. They are a very small subset: numbers, booleans, single characters. Most types in Rust are not Copy. And a type is not Copy because the compiler decides---it is because the programmer labels it. The compiler merely verifies that the label is compatible.''}

% ================================================================
\section{Heap Allocation with Box}
% ================================================================

All local variables live on the stack. This creates three situations where the stack is insufficient:

\begin{enumerate}
    \item \textbf{The data is too large} for the stack (stacks are typically \textasciitilde1~MB).
    \item \textbf{The data must outlive} the current function---when the function returns, the stack contracts and the data would evaporate.
    \item \textbf{The data size is not known at compile time}---the stack mechanism requires knowing exactly how many bytes to reserve at compilation.
\end{enumerate}

In these three situations, data must live on the \textbf{heap}. Rust provides \texttt{Box<T>} for this purpose.

\begin{memoryblock}
\textbf{Why the stack requires compile-time known sizes:} \\
The professor explained the technical reason behind the third constraint in detail: \emph{``The compiler, looking at how the function works, sees variable 1, variable 2, variable 3, and must immediately say: `variable 1 is located at stack pointer plus delta\_1, variable 2 at stack pointer plus delta\_2\ldots' with offsets calculated in the compilation phase. But if the data dimension were only known at runtime, these fixed offsets could not be stabilized.''} This is why dynamically-sized data \emph{must} go on the heap---the stack's addressing scheme is fundamentally based on fixed, compile-time-known offsets.
\end{memoryblock}

\begin{insightbox}
\textbf{Stack Overflow from Large Data} \\
The professor gave a concrete example of the ``data too large'' case: \emph{``A stack is typically one megabyte. If I need to place 20~GB---or even just 3~MB---I simply cannot put it on the stack. It would immediately get stuck. It would generate a stack overflow situation.''} This is why large data structures like big arrays or buffers must be heap-allocated using \texttt{Box}.
\end{insightbox}

\begin{lstlisting}[language=Rust]
let i = 5;                     // On the stack, no issues
let b = Box::new(27);          // 27 is allocated on the heap
println!("{}", *b);            // Dereference with * to access content
\end{lstlisting}

\texttt{Box} is a generic smart pointer: it holds a pointer to heap-allocated data and implements \texttt{Drop}. When the \texttt{Box} goes out of scope, it automatically frees the heap allocation.

\subsection{Box with Tuples: Detailed Walkthrough}

\begin{lstlisting}[language=Rust]
{
    let i = 5;                         // i on the stack (4 bytes)
    let mut b = Box::new((5_i32, 2_i32)); // Tuple on the heap
    (*b).1 = 7;                        // Modify second field via dereference
    println!("{:?}", b);               // Prints (5, 7)
} // b dropped: heap freed, then stack contracts
\end{lstlisting}

\begin{memoryblock}
\textbf{Step-by-step:}

\textbf{1.} \texttt{let i = 5;} --- 4 bytes on the stack for \texttt{i}. No special behavior.

\textbf{2.} \texttt{let mut b = Box::new((5, 2));} --- On the heap, a block of 8 bytes (two 4-byte \texttt{i32} values) is allocated, containing \texttt{(5, 2)}. On the stack, \texttt{b} occupies 8 bytes (a pointer to the heap block).

\textbf{3.} \texttt{(*b).1 = 7;} --- Dereference \texttt{b} with \texttt{*} to reach the heap, then access field \texttt{.1} (the second element of the tuple). Replace 2 with 7. The heap now contains \texttt{(5, 7)}.

\textbf{4.} The professor noted an important shorthand: the explicit \texttt{*} is needed only when you are not already accessing the content through the dot operator. In practice, \texttt{b.1 = 7;} also works because Rust automatically dereferences through \texttt{.} (auto-deref). The \texttt{println!} macro, for example, calls \texttt{b.debug()}, and the dot already implies dereferencing.

\textbf{5.} At the closing brace: \texttt{b} has drop instructions. The heap block (8 bytes) is freed first. Then the stack contracts to remove both \texttt{b} (8 bytes) and \texttt{i} (4 bytes).
\end{memoryblock}

\subsection{Returning a Box from a Function: Ownership Transfer}

\begin{lstlisting}[language=Rust]
fn make_box(a: i32) -> Box<(i32, i32)> {
    let r = Box::new((a, 1));
    return r; // Ownership of the Box transfers to the caller
}

fn main() {
    let b = make_box(5);
    println!("{:?}", b);     // Prints (5, 1)
} // b is dropped here: heap freed
\end{lstlisting}

\begin{memoryblock}
\textbf{What happens internally:}

\textbf{1.} In \texttt{main}, space is prepared on the stack for \texttt{b}, and then \texttt{make\_box(5)} is called.

\textbf{2.} Inside \texttt{make\_box}: the parameter \texttt{a} receives the value 5. The local variable \texttt{r} creates a \texttt{Box}: a heap block is allocated for the tuple \texttt{(5, 1)}, and \texttt{r} holds the pointer.

\textbf{3.} \texttt{return r;} --- This is a return assignment. Ownership of \texttt{r} is ceded to the caller's destination (\texttt{b}). Because ownership has transferred, \texttt{r} does \emph{not} execute its drop instructions---everything has passed to \texttt{b}.

\textbf{4.} \texttt{b} is now the owner of the heap block that was allocated by someone else (inside \texttt{make\_box}). The program proceeds; when \texttt{b} reaches the end of \texttt{main}, it frees the 8-byte heap block and the stack contracts.
\end{memoryblock}

% AI-IMAGE-REQUEST: Diagram showing make_box function call: (1) main's stack with space for b, (2) make_box's stack with parameter a=5 and local r pointing to heap (5,1), (3) return: arrow showing ownership transfer of heap pointer from r to b, (4) make_box stack frame removed but heap block persists, now owned by b in main.
\begin{figure}[!ht]
    \centering
    \includegraphics[width=0.85\textwidth]{images/ch04_box_return_ownership_transfer.png}
    \caption{Returning a \texttt{Box} from a function: the heap allocation survives the function's stack frame because ownership transfers to the caller.}
\end{figure}

\subsection{Box as a Specialized Pointer}

The professor placed \texttt{Box} within the broader context of Rust's pointer types during the recap lecture: \emph{``For the pointers, there are many types, each with different characteristics. The Box is a pointer that necessarily points to the heap and possesses its data. It's a pointer that, when it comes out of scope, guarantees the release of data. It's a pointer that is moved on assignment. A single variable is the owner of the Box. When that variable assigns it to someone else, the Box passes by property.''}

This makes \texttt{Box} one specific case in Rust's family of smart pointers. Others (like \texttt{Rc}, \texttt{Arc}, and references) have different ownership semantics, which we will explore in later chapters.

% ================================================================
\section{Borrowing: Access Without Possession}
% ================================================================

Constantly transferring ownership in and out of functions is tedious and unergonomic. The professor put it plainly: \emph{``When I want to pass data to a function, it's not always that I really want to `send it away' forever. If I just hand it over directly, Rust says `that's a transfer for me.'''}

To solve this, Rust introduces \textbf{references}---pointers that grant access \emph{without} ownership. This mechanism is called \textbf{borrowing}.

\begin{insightbox}
\textbf{The Library Book Analogy} \\
A reference is like borrowing a book from the library. You can read it, take notes from it, but you \emph{must return it}. You cannot burn it, sell it, or claim ownership. When the loan period ends, you have nothing to ``release''---the book continues to exist, controlled by its owner (the library). The reference allows access to a value \emph{without acquisition of the property}.
\end{insightbox}

Technically, a reference is still a pointer---an 8-byte number indicating a memory address. The difference from \texttt{Box} is entirely semantic: a reference carries \textbf{no ownership responsibility}. When a reference goes out of scope, nothing is released because the possession was someone else's and remains someone else's.

The professor positioned references in the broader taxonomy of pointers during the recap: \emph{``A reference is a pointer that does not possess the data---it is a loan. It means `I let you look inside, you can temporarily use it.' It is a pointer to a value that is possessed by another variable. When the reference comes out of scope, I don't have to do anything because the possession was of another and remains of the other.''}

\subsection{Immutable References (\texttt{\&T})}

An immutable reference grants \textbf{read-only} access to the underlying data. You can create as many immutable references simultaneously as you wish---multiple readers do not conflict.

\begin{lstlisting}[language=Rust]
fn print_string(s: &String) {
    println!("{}", s);     // Read access through reference
}                          // Reference dropped, nothing released

fn main() {
    let my_string = String::from("Hello");
    print_string(&my_string);  // Borrow, don't move!
    println!("{}", my_string); // Still valid! Ownership never left.
}
\end{lstlisting}

The \texttt{\&} symbol creates a reference: \texttt{\&my\_string} borrows the string without transferring ownership. Inside \texttt{print\_string}, the parameter \texttt{s} is a reference---it can look at the data but cannot modify or release it. When \texttt{print\_string} returns, the reference is simply dropped (no cleanup needed), and \texttt{my\_string} remains perfectly accessible.

\subsubsection{Owner Freeze During Immutable Borrows}

While immutable references exist, the underlying data is \emph{frozen}: the original owner is temporarily prohibited from modifying the data or transferring ownership. This guarantees that references always point to valid, stable values.

The professor elaborated on this freezing during the recap lecture: \emph{``While read-only references exist, the variable that possesses the value is also accessible---but only in read-only mode, even if we had potentially declared it mutable. It will return to full possession of its value only when the program reaches a line where none of the immutable references are still in use.''} This means the owner temporarily \emph{restricts its own power} because it has lent the data to others.

\subsection{Mutable References (\texttt{\&mut T})}

Sometimes you need to modify borrowed data in place. A mutable reference (\texttt{\&mut T}) grants this power, but under a strict condition: \textbf{exclusivity}.

\begin{lstlisting}[language=Rust]
fn add_world(s: &mut String) {
    s.push_str(" world");  // Modify through mutable reference
}

fn main() {
    let mut s = String::new();
    add_world(&mut s);        // Exclusive mutable borrow
    println!("{}", s);        // Prints " world"
}
\end{lstlisting}

\subsubsection{Owner Blocked During Mutable Borrows}

The professor detailed an even stricter rule for mutable references during the recap: \emph{``With mutable references, it is possible to get a mutable reference only from a mutable variable---otherwise I would give a permission that I don't have. And while there is a mutable reference, there cannot be any immutable references, nor is it possible to access the variable that actually possesses the value.''} The owner is completely locked out while the mutable borrow is active. \emph{``When I finally reach the line that goes beyond the last access to my mutable reference, the compiler unblocks the thing and leaves me to access it again.''}

This is more restrictive than immutable borrows: with immutable borrows, the owner can still \emph{read} the data. With mutable borrows, the owner cannot do \emph{anything} with the value until the borrow ends.

\subsection{The Cardinal Rule of Borrowing}

\begin{examprioritybox}
\textbf{The Readers--Writer Rule} \\
At any given time, you can have \textbf{either}:
\begin{itemize}
    \item \textbf{One mutable reference} (\texttt{\&mut T}) --- the single writer, or
    \item \textbf{Any number of immutable references} (\texttt{\&T}) --- multiple readers
\end{itemize}
\textbf{But never both simultaneously.}

The professor explained the reasoning: \emph{``If the writer is alone, no one gets hurt. If many are reading, no one gets hurt because nothing changes. But if while the writer is writing, someone reads, they could read something wrong---the writer wrote `ca' but is actually writing `casa.' The reader misinterprets. It must be certain that when I read, the data is stable.''}

This is the same principle found in databases (multiple readers / single writer locks) and concurrent programming. It prevents \emph{data races} at compile time.

The professor further explained during the recap: \emph{``The presence of the mutable reference guarantees that, for the time interval coinciding with the program interval in which the mutable reference is valid, no one can use another path to access that data. This is fundamental because if there is an in-progress transformation, readings made concurrently might read something wrong---intermediate states that tell me the data hasn't reached its stability yet.''}
\begin{itemize}
    \item \textbf{Expected Mastery Level:} Very High. You must identify violations of this rule in exam code snippets.
\end{itemize}
\end{examprioritybox}

% ================================================================
\section{The Borrow Checker and Lifetimes}
% ================================================================

The \emph{Borrow Checker} is the subsystem of the Rust compiler responsible for enforcing ownership and borrowing rules. It scrutinizes every code path---including all branches of \texttt{if} statements and all iterations of loops---to ensure that no reference outlives the data it points to.

\subsection{The Compiler Proves Theorems}

The professor described the Borrow Checker's work in formal terms during the recap lecture: \emph{``The compiler proves a theorem, basically. It proves many theorems. For any reference, it proves the theorem that, given that code, there is no path---no possible path within the various \texttt{if}, \texttt{while}, function calls, or such things---that allows contemporary access violating one of the rules we explained.''} This means the Borrow Checker is not performing simple pattern matching---it is conducting a formal proof over all possible execution paths.

\subsection{Lifetimes: The Concept}

Every value has a \textbf{lifetime}---the span of the program during which it is valid. A variable comes into existence when it is declared and ceases to exist when its scope ends. The professor drew a parallel to human life: \emph{``Before you are born, you don't exist. After you die, you no longer exist. In the middle is a piece of life. Programs revolve around time just the same.''}

Crucially, program execution is not simply linear. In a loop, line 23 may execute \emph{after} line 25 (because we loop back). So if line 25 gave something away (moved it), we cannot use it again when we return to line 23. We must follow the \textbf{logic of program execution}, tracing with our finger: first this, then this, then that, go back, this, this, go back...

The professor noted: \emph{``Programs break not because programmers are idiots. Simply because when complexity grows---nested ifs, loops, branching paths---the ability of the mind to track all ramifications becomes overwhelming. Some pieces will escape our attention. Rust's Borrow Checker does this tracking for us, following all paths, however complex, and verifying that in every possible execution path, the rules are never violated.''}

\subsection{What Lifetimes Really Are Internally}

The professor gave a more technical description of lifetimes during the recap lecture: \emph{``To be able to say whether there is concurrent access, the compiler must keep track, for each existing reference, of the set of lines in which that reference appears---in which it is valid, in which it is accessed. This set of lines generally takes the name of `lifetime.' It is information that lives inside the compiler's tables. The compiler uses it to do its math.''}

Lifetimes are fundamentally a compile-time bookkeeping mechanism. They have no runtime representation whatsoever. The compiler uses them to prove that the \textbf{lifetime of a reference is completely contained within the lifetime of the value it points to}---and if it cannot prove this, compilation fails.

\subsection{The Core Lifetime Rule}

The lifetime of a reference must be \textbf{entirely contained within} the lifetime of the value it points to. The professor used a body analogy: \emph{``If your liver is destined to be destroyed tomorrow, how long can you live? At most until tomorrow---because an essential part of you ceases to exist. You might die sooner (a brick falls on your head), but never later.''}

This propagation works recursively: if a reference is inserted into a larger structure, the lifetime of the entire structure is constrained to the lifetime of the reference. And if that structure is part of an even larger one, the constraint propagates upward.

\subsection{Dangling Reference Example}

\begin{lstlisting}[language=Rust]
let r;
{
    let x = 1;
    r = &x;       // r borrows x
}                  // x is dropped here!
// println!("{}", r); // ERROR: x does not live long enough
\end{lstlisting}

\begin{memoryblock}
\textbf{What happens internally:}

\textbf{1.} \texttt{r} is declared on the stack (8 bytes for a pointer), but not yet initialized.

\textbf{2.} A new block opens. \texttt{x = 1} is placed on the stack (4 bytes).

\textbf{3.} \texttt{r = \&x;} --- \texttt{r} now points to the stack location of \texttt{x}.

\textbf{4.} The closing brace contracts the stack, removing \texttt{x}. But \texttt{r} still holds the old address---now pointing to deallocated stack memory. If an interrupt fires, it could overwrite that memory. If we read \texttt{*r}, we might get garbage.

\textbf{5.} In C, this compiles and ``works'' 49 out of 50 times. The 50th time, an interrupt happens at the wrong moment and you get corrupted data. \textbf{In Rust, the compiler catches this at compile time} and says: \emph{``\texttt{x} does not live long enough---it is dropped at the closing brace while still borrowed by \texttt{r}.''}
\end{memoryblock}

The professor used the rental analogy: \emph{``It's like when your lease on a storage unit ends and you return the keys. But you made a copy. The landlord rents it to someone else. You walk in with your copy of the key---but the fridge you expected to find is gone, replaced by someone else's furniture. That's a dangling pointer.''}

\subsubsection{Dangling Pointers in C: A Detailed Walkthrough}

The professor gave a concrete walkthrough of a dangling pointer in C during the recap lecture: \emph{``A dangling pointer is something I create in C with \texttt{malloc}. I do a 10-byte \texttt{malloc}. I get the pointer: \texttt{0x7B5F4N}. I save it in a variable. At this moment, it exists. I use \texttt{*ptr} to access the object. Then at some point I \texttt{free} that pointer and return the memory. But I don't cancel the pointer---it remains written \texttt{0x7B5F4N} somewhere. In the next few lines, if I hadn't noticed that I had freed it and I still use it, the compiler would let me do it. It's a shame that my house is no longer there. It became someone else's house---that piece of memory is now used by something else. You go inside and do things. That's a dangling pointer---a pointer that is used outside of the time in which it is valid.''}

\emph{``Since C doesn't track the lifetime, the compiler doesn't know how long a value lasts. That awareness is entirely the programmer's responsibility. If the programmer does well, the problem doesn't arise. The real problem is that the programmer often doesn't do well---not because they are bad, but because the details are many, and something will escape.''}

\subsection{Lifetime Propagation into Containers}

\begin{lstlisting}[language=Rust]
let mut v = Vec::new();
{
    let data = 1;
    v.push(&data);  // Reference to data stored in vector
}                    // data is dropped!
// println!("{:?}", v); // ERROR: v contains a dangling reference
\end{lstlisting}

The vector \texttt{v} contains a reference to \texttt{data}. When \texttt{data} is dropped, that reference becomes dangling. Therefore the \emph{entire} vector becomes invalid. The compiler reports: \emph{``You took a reference to \texttt{data} at line 4, but \texttt{data} does not live long enough---it was dropped at line 5 while still borrowed.''}

The professor generalized: \emph{``If I had put five references in the vector, the vector's lifetime would automatically be the shortest among all the things I put in it. The lifetime of the weakest part propagates to the container.''}

The professor further explained during the recap: \emph{``If this reference becomes part of a wider structure, the constraint is automatically propagated to the wider structure. And if that structure is still part of another thing, it goes up recursively. The same for functions: if a function receives a reference as a parameter, the lifetime of that reference becomes part of the function's signature. The function is usable only while that reference exists. The time of life of the reference conditions the time of life of the function.''} When you choose a specific reference to pass to a function, you are conditioning the usability of the function itself.

\subsection{Lifetime Annotations}

In most cases, the Borrow Checker can deduce lifetimes automatically (\textbf{lifetime elision}). However, when functions receive multiple references and return a reference, the compiler cannot always determine which input lifetime the output depends on. In these cases, \textbf{explicit lifetime annotations} are required:

\begin{lstlisting}[language=Rust]
fn longest<'a>(s1: &'a str, s2: &'a str) -> &'a str {
    if s1.len() > s2.len() { s1 } else { s2 }
}
\end{lstlisting}

The annotation \texttt{'a} is a ``lifetime parameter''---it tells the compiler that the returned reference lives \emph{as long as both inputs}. The professor explained: \emph{``I could have a function `max' that takes references to two numbers and returns a reference to the larger one. How long does the result last? It depends on which input was selected---and the compiler can't know that statically. So we must tell it: the result lasts as long as both inputs.''}

\subsubsection{Lifetime Annotations as Purely Symbolic Names}

The professor gave a particularly illuminating explanation of what lifetime annotations actually represent during the recap lecture:

\emph{``We don't give a concrete representation to our lifetime. We're not going to say `you go from January 25th to January 27th'---that doesn't make any sense. We give the references' lifetimes a purely symbolic name: \texttt{'a}, \texttt{'b}, \texttt{'c}, where \texttt{a}, \texttt{b}, and \texttt{c}, in themselves, have no meaning, other than that within the same function, all things marked with \texttt{'a} mean the same thing, and they are different from things marked \texttt{'b}.''}

\emph{``So the lifetime criterion is only a way to say `this is the same as that' or `this is different from that.' The name is purely arbitrary.''}

This leads to an important consequence for shared lifetimes: \emph{``If one reference could go from line 1 to line 11, and the other from line 3 to line 13, giving both the same lifetime \texttt{'a}, automatically---even though one could individually be valid for a greater number of lines---because I'm forcing them to use the same lifetime, I'm tightening the usage to only the intersection. Trying to access one or the other outside the intersection will cause a compilation error.''}

\begin{exambox}
\textbf{Lifetime Intersection Rule} \\
When two references share the same lifetime parameter \texttt{'a}, their valid usage window becomes the \textbf{intersection} of both lifetimes. If reference A is valid from line 3 to line 13 and reference B is valid from line 1 to line 11, then both are usable only from line 3 to line 11. Accesses outside this intersection are compile errors.
\end{exambox}

Lifetime rules for functions:
\begin{itemize}
    \item If there is a \emph{single} reference parameter, the compiler automatically assumes the output lifetime equals the input lifetime.
    \item If there are \emph{multiple independent} references, the compiler may require explicit annotations.
    \item The presence of a lifetime parameter in a function signature conditions the \emph{usability} of the function: you can only call it while all the referenced data is alive.
    \item If a function receives more references, it could be necessary to indicate explicitly whether those references must last at the same time or if they are independent.
\end{itemize}

\subsection{Lifetimes in Structs}

When a struct holds a reference, the struct's lifetime is bounded by the reference's lifetime:

\begin{lstlisting}[language=Rust]
struct User<'a> {
    id: u32,
    name: &'a str,  // Reference with lifetime 'a
}
\end{lstlisting}

The \texttt{User} struct can only live as long as the \texttt{name} reference is valid. The \texttt{id} field is owned by the struct (it is just a number), so it poses no constraint. But \texttt{name} is borrowed---and this borrowed piece makes the entire container perishable. If \texttt{User} is placed inside yet another struct, the lifetime constraint propagates recursively upward.

\subsection{The \texttt{'static} Lifetime}

The special lifetime \texttt{'static} means the data lives for the \emph{entire duration of the program}. String literals are \texttt{\&'static str} because they are stored in the read-only constants segment of the binary, which is never freed:

\begin{lstlisting}[language=Rust]
let s: &'static str = "Hello"; // Lives forever
\end{lstlisting}

This is different from taking a reference to a \texttt{String} created at runtime, which would have a limited lifetime.

The professor elaborated during the recap: \emph{``When we write a string literal like \texttt{"ciao"}, it automatically becomes a compile-time constant that exists for the entire duration of the program because it is written directly into the executable, somewhere, and it will be there the entire time. So we can, at any moment in our program's life, refer to that particular string---but exactly that one, which will be at some fixed address. \texttt{\&'static str} denotes a sequence of characters that exists for the duration of the program.''} It is fundamentally different from \texttt{\&str} obtained by borrowing a \texttt{String}: the latter is \texttt{\&'a str} where \texttt{'a} is the lifetime of the \texttt{String}, which ``lasts little'' by comparison.

\subsection{All Checks Are Compile-Time}

\begin{insightbox}
\textbf{Zero Runtime Cost} \\
The professor emphasized: \emph{``All of this is done at compile time. There is no trace of it at runtime. It doesn't affect execution time. We don't pay for it. It's not like Java's null pointer check, where every time Java has a pointer it asks `is this null?' and burns four CPU cycles just for that---and pointers are everywhere. Here, the control is done a priori. References can never be null: if I can only create a reference from data that exists, null references are ontologically impossible.''}
\end{insightbox}

The professor further explained the null-impossibility argument: \emph{``\texttt{r = \&a}. If \texttt{a} is not there, I can't take the reference. So the null reference doesn't exist. If there are references, it means the data existed. But that reference is accessible only while the data exists.''}

% ================================================================
\section{Slices: References to Sequences}
% ================================================================

In C, a pointer to an array decays into a pointer to the first element, losing all dimension information. This is the root cause of buffer overflows and the notorious complexity of C string handling.

Rust addresses this with \textbf{slices}: a reference to a contiguous sequence of elements, represented as a \textbf{fat pointer}.

\subsection{Fat Pointers}

A regular reference (\texttt{\&T}) occupies 8 bytes---just a memory address. A slice (\texttt{\&[T]}) occupies \textbf{16 bytes} (128 bits):

\begin{itemize}
    \item First 8 bytes: pointer to the first element
    \item Second 8 bytes: the number of elements (length)
\end{itemize}

This explicit length allows Rust to perform \textbf{bounds checking at runtime}: if you attempt to access an index beyond the length, the program will \textbf{panic} (a controlled crash) rather than silently corrupting memory.

\subsection{Creating Slices}

\begin{lstlisting}[language=Rust]
let a = [3, 4, 5, 6, 7]; // Fixed-size array on the stack
let s: &[i32] = &a[1..3]; // Slice: elements at index 1 and 2
// s points to [4, 5], length = 2
println!("{:?}", s);       // Prints [4, 5]
\end{lstlisting}

The notation \texttt{\&a[1..3]} creates a reference to elements from index 1 (inclusive) to index 3 (exclusive): values 4 and 5. The slice \texttt{s} is a fat pointer containing the address of element at index 1 and the count 2.

\begin{warningbox}
\textbf{Slices and Panic} \\
Accessing a slice out of bounds triggers a \textbf{panic}---the program terminates in a controlled manner. This is dramatically safer than C's silent buffer overflow. However, you should still write defensive code: check the slice length before indexing, use iterators when possible, and avoid raw index access when the bounds are uncertain. The slice index type is \texttt{usize} (unsigned), so negative indices are syntactically impossible.
\end{warningbox}

The professor noted: \emph{``The slice encapsulates the length, so I am able to verify that we're not going beyond our responsibility. The only way to guarantee bounds safety is at runtime (since indices depend on runtime values), and slices make this possible by carrying their length with them.''}

Slices can be taken from arrays, vectors, strings (\texttt{\&str} is a string slice), and even from \texttt{Box} allocations:

\begin{lstlisting}[language=Rust]
let v = vec![10, 20, 30, 40, 50];
let slice: &[i32] = &v[1..4]; // [20, 30, 40]
println!("{:?}", slice);
\end{lstlisting}

Since a slice is a reference, all borrowing rules apply: while the slice exists, the vector is borrowed and cannot be modified.

% AI-IMAGE-REQUEST: Diagram showing an array [3,4,5,6,7] on the stack, with a fat pointer (slice) consisting of two 8-byte fields: a pointer to element at index 1 (value 4) and length=2. Show the slice "window" highlighting elements [4,5] within the array.
\begin{figure}[!ht]
    \centering
    \includegraphics[width=0.8\textwidth]{images/ch04_slice_fat_pointer.png}
    \caption{A slice as a fat pointer: the pointer addresses the first included element, and the length field enables runtime bounds checking.}
\end{figure}

% ================================================================
\section{Summary: Box vs.\ Reference --- A Comparison}
% ================================================================

The professor presented a clear comparison between the two main pointer types:

\begin{center}
\begin{tabular}{|l|c|c|}
\hline
\textbf{Property} & \textbf{\texttt{Box<T>}} & \textbf{\texttt{\&T} / \texttt{\&mut T}} \\
\hline
Ownership & \textbf{Owns} the data & \textbf{Borrows} the data \\
Points to & Always the heap & Stack or heap \\
On drop & Frees the heap allocation & Does nothing \\
Duplicable? & No (moves on assignment) & Yes (\texttt{\&T} can be copied freely) \\
Mutable access? & If the Box is mutable & Only with \texttt{\&mut T} (exclusive) \\
Can be null? & No & No \\
Binary representation & 8 bytes (pointer) & 8 bytes (pointer) \\
\hline
\end{tabular}
\end{center}

The professor noted: \emph{``Both Box and Ref are pointers. If I look at them in binary representation, they are indistinguishable---8 bytes with a number. The difference is entirely in semantics. Box is responsible for release. Ref is not---the data exists before the reference and will continue after it.''}

The professor also framed this comparison during the recap: \emph{``I can copy a reference---there can be many in read-only mode, or one in write mode. I can't copy a Box---I can move it, but there cannot exist two Boxes that point to the same data.''}

Key distinctions from C pointers:
\begin{itemize}
    \item Neither \texttt{Box} nor \texttt{\&T}/\texttt{\&mut T} can be null. There is no possibility of null pointer dereferencing.
    \item C pointers allow aliasing (multiple pointers to same data) without any control, creating validity and double-free problems. With \texttt{\&T}, many can read simultaneously but the double-free problem cannot arise because references don't release. With \texttt{\&mut T}, exclusivity prevents data races. With \texttt{Box}, there is always a unique owner.
    \item In C, a pointer might represent a single value or an array (the type doesn't distinguish). Rust uses slices (\texttt{\&[T]}) as fat pointers to represent arrays, encoding the length explicitly.
\end{itemize}

\begin{insightbox}
\textbf{The Great Ambiguity of C, Resolved} \\
The professor concluded the pointer comparison with a summarizing observation: \emph{``Despite all this complexity, it is actually to protect us. The effort we have to make to squeeze awareness, to disambiguate situations---the great ambiguity of the C pointers---here we fight with different pointers, each of which has powers and limits. This allows the compiler to do a control that it would not otherwise be able to do, and to generate only the necessary operations, giving us access to what we need.''} The ``complexity'' of Rust's pointer system is the \emph{solution} to C's dangerous simplicity.
\end{insightbox}

% ================================================================
\section{Why Ownership Matters: Performance Without GC}
% ================================================================

The professor concluded the lecture by connecting ownership to performance guarantees. Rust's ownership system provides all benefits of a garbage collector \emph{without one}:

\begin{itemize}
    \item \textbf{No non-deterministic pauses:} Programs have performance homogeneous with respect to the same input. Garbage-collected programs depend on factors outside your control, introducing non-deterministic pauses.
    \item \textbf{Minimal memory usage:} GC languages tend to hold memory until a collection cycle runs (``spray, spray, spray, and eventually I'll recover something''). Rust releases memory deterministically the moment it's no longer needed.
    \item \textbf{Cloud cost implications:} The professor noted that in cloud environments where memory usage is billed, writing programs that use a tenth of the memory means paying significantly lower prices.
    \item \textbf{Resource cleanup beyond memory:} Ownership and \texttt{Drop} guarantee correct release of file handles, sockets, hardware devices---not just memory.
\end{itemize}

\begin{insightbox}
\textbf{Reading the Compiler Messages} \\
The professor urged students to \emph{read} the Rust compiler's error messages carefully: \emph{``Students arrive, write their code, run \texttt{cargo build}, it breaks. `Fix it.' `I don't know how.' `How can you not know? It's written right there!' The compiler explains in detail what's wrong and why. It's verbose, yes, but extremely analytical. It says exactly: `Look, the problem is here.' The solution is ours to find, but the diagnosis is precise. You must not be intimidated by the number of lines---force yourself to read them.''}

The professor added: \emph{``The compiler can't tell us how to adjust it, but it explains in detail why something doesn't compile. Either we eliminate the assignment, or we eliminate the print line. The solution is for us to decide---but the problem is precisely identified.''}
\end{insightbox}

% ================================================================
\section{Ownership Recap: The Professor's Summary}
% ================================================================

At the beginning of the second lecture session, the professor provided a structured recap of all ownership concepts. This summary is valuable as a mental checklist:

\begin{exambox}
\textbf{The Ownership System in Summary} \\
\begin{enumerate}
    \item Every value is possessed by one and only one variable. That variable is responsible for its release.
    \item The \textbf{type} of the value indicates what must be done at release time. Simple types require only the elimination of the variable (stack contraction). More articulated types require specialized instructions (\texttt{Drop} trait).
    \item When I assign a variable to another, pass a variable as a parameter, or return a variable as a result---in all three cases, the \textbf{possession transits} from the source to the destination through the operation called \emph{move}. The source becomes inaccessible.
    \item Some particular types are marked as \texttt{Copy}. This means: since there are no particular recycling instructions and the meaning of what is written inside has no special implications, the information remains accessible after a copy.
    \item When a pointer \textbf{possesses} the data it points to, use the \texttt{Box} type. Box is useful for variable-length data, very large data, or data that must outlive its function.
    \item When I need to \textbf{access without owning}, use \textbf{references} (\texttt{\&T} or \texttt{\&mut T}). References are pointers without possession---loans.
    \item The Borrow Checker is the compiler component that manages references. It enforces that \texttt{\&T} and \texttt{\&mut T} are mutually exclusive: many \texttt{\&T}, or one \texttt{\&mut T}, never both.
    \item While mutable references exist, even the original owner cannot access the data.
    \item All checks are done at compile time with zero runtime cost.
\end{enumerate}
\end{exambox}

% ================================================================
\section{Comprehensive Ownership Tracing: Exam-Style Example}
% ================================================================

\begin{exambox}
\textbf{Mental Model for the Exam} \\
When analyzing Rust code, mentally track who \emph{owns} each value and trace the \emph{lifetimes} of any borrows. The professor walked through this exact pattern:

\begin{lstlisting}[language=Rust]
let s1 = String::from("Hello");  // s1 owns heap data
let s2 = s1;                      // MOVE: s1 -> s2
// s1 is now inaccessible
let num1 = 3;                     // num1 on stack (Copy type)
let num2 = num1;                  // COPY: num1 still valid
// Both num1 and num2 are independently usable
\end{lstlisting}

The stack trace:
\begin{enumerate}
    \item \texttt{s1} allocated (24 bytes), heap block for ``Hello''
    \item \texttt{s2} allocated (24 bytes), receives bitwise copy of \texttt{s1}'s stack data. \texttt{s1} marked inaccessible.
    \item \texttt{num1} allocated (4 bytes), value = 3
    \item \texttt{num2} allocated (4 bytes), value = 3 (copy, not move---\texttt{num1} remains valid)
\end{enumerate}

At end of scope:
\begin{itemize}
    \item \texttt{num2}: no special drop, stack contracts
    \item \texttt{num1}: no special drop, stack contracts
    \item \texttt{s2}: Drop trait invoked $\rightarrow$ heap freed, then stack contracts
    \item \texttt{s1}: Drop \emph{omitted} (no longer owns anything), stack contracts
\end{itemize}

The professor noted an important observation: \emph{``\texttt{num1} can increase freely and \texttt{num2} goes for its road. They are two independent copies and there are no particular links between them.''}
\end{exambox}

% ================================================================
\section*{Exam Relevance and Recurrence}
% ================================================================

\textbf{Score: Very High --- Tested in EVERY exam} \\
\textbf{Notes:} This topic forms the absolute core of the exam. You will consistently find questions asking you to:
\begin{itemize}
    \item Trace ownership through a code snippet, identifying which variables own which values at each point
    \item Identify compilation errors caused by moves, dangling references, or borrowing violations
    \item Explain the difference between a bitwise copy (Move) and a deep copy (\texttt{Clone})
    \item Explain why \texttt{Copy} and \texttt{Drop} are mutually exclusive
    \item Determine the lifetime of references and identify dangling pointer situations
    \item Distinguish between shadowing (\texttt{let x = ...}) and reassignment (\texttt{x = ...})
    \item Explain how the owner is frozen/blocked during borrows (immutable: read-only; mutable: completely blocked)
    \item Describe why lifetime annotations use symbolic names (\texttt{'a}, \texttt{'b}) and what shared lifetime parameters imply (intersection)
\end{itemize}

% ================================================================
\section{Domande ed Esercizi da Esami Passati}
% ================================================================

\begin{itemize}
    \item \begin{exambox}
    \textbf{Domanda:} Si definiscano i concetti di Dangling Pointer, Memory Leakage e Wild Pointer, facendo esempi concreti, usando dello pseudocodice, che possono generare questi fenomeni.
    
    \textbf{Risposta:} In base al capitolo, un \emph{dangling pointer} è un puntatore a un'area di memoria non più valida o deallocata. In C, si verifica ad esempio se si usa \texttt{malloc}, si salva il puntatore, si esegue la \texttt{free} per restituire la memoria, ma non si cancella il puntatore originale, continuando ad usarlo (come avere la copia della chiave di un magazzino già affittato ad altri). Il \emph{memory leakage} avviene invece quando la memoria allocata non viene mai rilasciata (nessun proprietario chiama le istruzioni di rilascio o \texttt{Drop}). I \emph{wild pointer} non sono direttamente citati nel testo, ma rientrano nella problematica generale dell'ambiguità dei puntatori privi di validità garantita a compile-time. Rust previene tutto ciò assegnando ad ogni dato un solo \emph{owner} e assicurando che la \emph{lifetime} dei puntatori (reference) non superi mai quella del valore originario.
    \end{exambox}

    \item \begin{exambox}
    \textbf{Domanda:} Completa questa funzione in modo che compili correttamente, se ritieni che ci sia un errore di compilazione (descrivendolo), altrimenti dichiara il motivo per cui la funzione è scritta correttamente:
\begin{lstlisting}[language=Rust]
fn get_prefix(s: &str, len: usize) -> &str {
    &s
\end{lstlisting}

    \textbf{Risposta:} La funzione è incompleta. Per compilare correttamente e restituire un prefisso (ovvero una \emph{string slice}), deve essere completata usando la sintassi delle \emph{slice}:
\begin{lstlisting}[language=Rust]
fn get_prefix(s: &str, len: usize) -> &str {
    &s[..len]
}
\end{lstlisting}
    Non ci sono errori di compilazione legati alle \emph{lifetime} perché il compilatore applica la regola di \emph{lifetime elision}: essendoci un solo parametro di tipo reference in ingresso (\texttt{s: \&str}), assume automaticamente che la lifetime del valore restituito sia identica alla lifetime dell'input.
    \end{exambox}

    \item \begin{exambox}
    \textbf{Domanda:} Si consideri il seguente frammento di programma Rust:
\begin{lstlisting}[language=Rust]
1 fn main() {
2 let mut s = String::from("esame");
3 let r1 = &s;
\end{lstlisting}

    \textbf{Risposta:} In questo frammento, \texttt{s} è il proprietario (\emph{owner}) di una stringa allocata nell'heap e \texttt{r1} è una \emph{immutable reference} (prestito o \emph{borrow} in sola lettura) a \texttt{s}. Secondo le regole del \emph{Borrow Checker}, fintanto che \texttt{r1} è in uso, il proprietario originale \texttt{s} viene "congelato" (\emph{owner freeze}): non può modificare la stringa, né trasferirne la proprietà, per garantire che \texttt{r1} continui a puntare a dati validi e stabili. Questo previene la possibilità di avere modifiche impreviste mentre vi sono lettori attivi (regola dei \emph{Readers-Writer}).
    \end{exambox}

    \item \begin{exambox}
    \textbf{Domanda:} Si descriva la gestione della memoria in Rust e si spieghi come vengono evitati i problemi di sicurezza comuni come le violazioni di accesso o la presenza di puntatori nulli.
    
    \textbf{Risposta:} Rust gestisce la memoria tramite il concetto di \emph{Ownership} valutato a tempo di compilazione (zero runtime cost). Ogni valore ha un solo \emph{owner} e viene rilasciato automaticamente all'uscita dello scope, prevenendo \emph{double free} o \emph{memory leak}. I prestiti (\emph{borrowing}) evitano copie non necessarie tramite reference (\texttt{\&T} o \texttt{\&mut T}); le violazioni di accesso o le data race sono impedite assicurando la mutua esclusione tra \texttt{\&T} (molti lettori) e \texttt{\&mut T} (un solo scrittore esclusivo). I puntatori nulli non esistono perché ogni reference deve essere sempre valida (\emph{lifetime} verificata staticamente), mentre per identificare i limiti degli array Rust usa i \emph{fat pointer} (\emph{slice}), garantendo l'assenza di buffer overflow tramite runtime bounds checking.
    \end{exambox}

    \item \begin{exambox}
    \textbf{Domanda:} Si descriva il concetto di "ownership" in Rust e come contribuisca a prevenire errori come le race condition e i dangling pointer rispetto ad altri linguaggi come C++.
    
    \textbf{Risposta:} L'\emph{ownership} stabilisce che ogni valore ha uno e un solo proprietario, responsabile della sua deallocazione (\texttt{Drop}). Questo previene problemi come \emph{double free} perché solo l'owner rilascia la risorsa. Nel concedere accessi temporanei (\emph{borrowing}), il compilatore verifica che una reference non viva più a lungo del suo proprietario (\emph{lifetime core rule}), eliminando alla radice i \emph{dangling pointer}. Le \emph{race condition} sono prevenute dalla rigida regola dei Readers-Writer: a compile-time non è possibile avere referenze mutabili (\texttt{\&mut T}) contemporaneamente ad altre referenze alla stessa memoria, garantendo stabilità dei dati.
    \end{exambox}

    \item \begin{exambox}
    \textbf{Domanda:} Si spieghino le differenze tra costruzione di copia e costruzione per movimento nel linguaggio C++11 in termini di comportamento e di prestazioni, e si descriva un caso d'uso nel quale il movimento risulta vantaggioso rispetto alla copia.
    
    \textbf{Risposta:} Applicando i principi descritti nel capitolo al concetto di "movimento", una \emph{move} in Rust è una copia a basso costo (\emph{bitwise copy} o \texttt{memcpy} dello stack, tipicamente 24 byte per una Stringa) che contemporaneamente invalida l'accesso alla variabile sorgente, evitando di dover clonare o gestire i dati nell'heap. La costruzione di \emph{copia} completa (in Rust rappresentata da \texttt{Clone}) genera invece una nuova allocazione nell'heap, implicando un costo di prestazione proporzionale ai dati (es. ricopiare l'intera "Divina Commedia"). Il movimento è estremamente vantaggioso ovunque sia necessario trasferire il possesso di un oggetto "pesante" (es. \texttt{String} o array dinamico) senza sostenere l'onere in termini di tempo e memoria richiesto dalla duplicazione dei dati heap associati.
    \end{exambox}
\end{itemize}
